\documentclass[journal,10pt,twocolumn]{IEEEtran}

% ---- Packages ----
\usepackage[utf8]{inputenc}
\usepackage[english]{babel}
\usepackage{amsmath, amssymb}
\usepackage{bm}
\usepackage{graphicx}
\usepackage{booktabs}
\usepackage{multirow}
\usepackage{algorithm}
\usepackage{algpseudocode}
\usepackage{hyperref}
\usepackage{cite}

\hypersetup{colorlinks=false}

% ============================================================
\title{Distributed Multi-Drone ARTVA Search \& Rescue\\ in GPS-denied scenarios}

\author{Andrea~Renelli,~\IEEEmembership{Student,~University~of~Trento}
        and~Luca~Zanolini,~\IEEEmembership{Student,~University~of~Trento}%
\thanks{University of Trento, Italy.
  \texttt{\{andrea.renelli, luca.zanolini-1\}@studenti.unitn.it}.}
\thanks{This report is submitted as the final project for the course
  ``Intelligent Distributed Systems''.}}

% ============================================================
\begin{document}

\maketitle

% ============================================================
\begin{abstract}
In avalanche rescue, survival probability drops sharply within the first 15 minutes of burial, making rapid victim localization critical. Deploying a coordinated drone fleet can drastically reduce search time compared to manual ARTVA sweeps, but requires solving path planning, source localization, and distributed coordination under strict real-time constraints.
This report presents a fully distributed multi-drone system for ARTVA source localization with no central coordinator: each agent maintains only local state and communicates within a limited radio range. Each drone operates as a Finite State Machine cycling through \texttt{SEARCH}, \texttt{TRACK}, \texttt{STOP}, and \texttt{SUPPORT} states, with motion governed throughout by a Model Predictive Controller. In GPS-denied conditions, state estimation relies on IMCDL (Interim Master Decentralized Cooperative Localization), combining inter-drone UWB ranging and LiDAR altitude to limit odometric drift.
Signal intensity thresholds drive state transitions: drones explore freely in \texttt{SEARCH}, converge toward the source via an Extremum Seeking controller in \texttt{TRACK}, then halt and recruit two support drones for triangulation via min-consensus. Source estimation uses Distributed Cooperative Gradient Descent with ATC diffusion.
A parametric sweep of [XXXX] experiments across five alpine terrain patches, three fleet sizes ($N \in \{3,4,5\}$), three noise levels, and four communication radii achieves an overall mission success rate of [XX.X]\,\% with a median search time of [XX]\,s.
\end{abstract}

\begin{IEEEkeywords}
Multi-drone systems, ARTVA search and rescue, Model Predictive Control,
Extremum Seeking, IMDCL cooperative localization, distributed gradient descent,
avalanche rescue.
\end{IEEEkeywords}

% ============================================================
\section{Introduction}
Avalanche rescue is a race against time: survival probability drops
dramatically after the first 15 minutes of burial~\cite{artva2020}.
Traditional ARTVA searches rely on rescuers sweeping the terrain on foot,
a process that is both slow and physically demanding in avalanche conditions.
Autonomous aerial vehicles offer a promising alternative: a coordinated
drone fleet can cover the search area far faster than a human sweep. In a realistic deployment, first
responders arriving on site release the fleet, which autonomously scans
the terrain, localises the buried victim, and marks the estimated position
for excavation, either by landing on the target or by releasing on the ground a light emitter visible through the snowpack. The challenge lies in coordinating multiple drones, each operating on partial, noisy measurements and peer-to-peer communication alone, to converge on the source without a central coordinator, a global map, or external positioning infrastructure.

Centralised architectures are fragile in avalanche scenarios: a crashed coordinator
can abort the entire mission.  Fully distributed systems, where each drone operates
on local measurements and peer-to-peer messages, are inherently resilient but must
resolve four tightly coupled sub-problems.
First, each drone must track a sequence of mission waypoints subject to actuator
constraints, using only its locally estimated state.
Second, in the absence of GPS, position estimates drift without an
external absolute reference, forcing the fleet to maintain accurate
self-localisation through inter-drone ranging and on-board sensing alone.
Third, the fleet must cooperatively fit a nonlinear source model to noisy,
spatially distributed ARTVA measurements.
Fourth, once a drone detects a strong enough signal it must recruit the nearest
reachable peers for support through purely local communication.

Related work on distributed ARTVA search~\cite{artva2020} motivates the aerial
platform choice.  The IMDCL framework of Kia et al.~\cite{kia2016cooperative}
provides cross-covariance-aware decentralised filtering without a global state
vector.  Olfati-Saber and Murray~\cite{olfati2004consensus} and Chen and
Sayed~\cite{chen2012diffusion} establish the consensus and diffusion backbone of
the source-estimation approach.  The ES formulation follows Azzollini et
al.~\cite{azzollini2021es}, whose SITL evaluation demonstrates convergence from
$\sim50\,$m away under conditions comparable to this deployment.

\subsection*{Contributions and paper organisation}

The main contributions of this work are: (i)~a complete simulation of a
distributed multi-drone ARTVA search system on real TINItaly (\textit{INGV}) terrain,
integrating MPC, IMDCL, ES, and DCGD without any central coordinator;
(ii)~a mission state machine coupling four distinct drone behaviours—\textsc{Search},
\textsc{Track}, \textsc{Stop}, \textsc{Support}—with automatic noise calibration
and threshold scaling via average consensus; (iii)~a min-consensus partner
selection protocol that identifies the $P=2$ nearest reachable drones for
triangulation support in at most $K_{\max} = 10$ communication rounds; and
(iv)~a systematic parametric sweep of [XXXX] experiments demonstrating [XX.X]\,\%
overall mission success and isolating the impact of fleet size, noise, communication
radius, and burial depth.

Section~\ref{sec:models} describes the drone dynamics and signal models.
Section~\ref{sec:solution} details the four core algorithms.
Section~\ref{sec:impl} covers implementation aspects.
Section~\ref{sec:results} presents the experimental sweep and discussion.
Section~\ref{sec:conclusions} concludes.

% ============================================================
\section{Adopted Models}
\label{sec:models}

\subsection{Communication Model}

At each simulation step $t$, drones communicate over an undirected
graph $\mathcal{G}(t) = (\mathcal{V}, \mathcal{E}(t))$ where
$\mathcal{V} = \{1,\ldots,N\}$ and an edge $(i,j) \in \mathcal{E}(t)$
exists iff $\|\mathbf{p}_i(t) - \mathbf{p}_j(t)\| \leq R_c$.
The nominal $R_c = 80\,$m models open-alpine UWB conditions and it is swept in the
experiments (Section~\ref{sec:sweep}).
Communication is symmetric and range-limited; no agent has global topology
information.

\subsection{Drone Dynamics}
Since the main purpose of this project is to return a proof of concept for the whole framework, each drone is modelled as a \textbf{3D point mass} controlled in acceleration. This abstraction is justified by two observations: first, the ARTVA signal depends only on position and not on orientation, so attitude dynamics do not affect the measurement model; second, the point-mass layer can be treated as a black box hiding a
lower-level attitude controller that, given the desired acceleration
computed by the high-level controller (in this case an MPC), allocates motor thrusts and stabilises the vehicle
orientation, as is standard in quadrotor cascaded control architectures.
States and inputs are:
$x_i = [p_x,\,p_y,\,p_z,\,v_x,\,v_y,\,v_z]^\top$ and $u_i = [a_x,\,a_y,\,a_z]^\top$.
The continuous dynamics follow the standard point-mass model
$\dot{x} = A_c\,x + B_c\,u$.
The Zero-Order Hold (\textbf{ZOH}) discretisation (step $dt = 0.1\,$s)
assumes the input $u$ is held constant over each interval
$[k\,dt,\,(k{+}1)\,dt]$; since $A_c$ and $B_c$ encode simple Newtonian
kinematics, the matrix exponential resolves in closed form and the
discretisation is exact, giving:
\begin{equation}
  x_i(k+1) = F\,x_i(k) + G\,u_i(k) + G\,\eta_i(k),
  \label{eq:dynamics}
\end{equation}
where $\eta_i \sim \mathcal{N}(0, Q)$ is isotropic process noise with
$\sigma_{acc} = 0.05\,$m/s$^2$ (see Section~\ref{sec:sweep} for the sweep), and:
\[
  F = \begin{bmatrix} I_3 & dt\,I_3 \\ 0_3 & I_3 \end{bmatrix},
  \quad
  G = \begin{bmatrix} \tfrac{1}{2}dt^2\,I_3 \\ dt\,I_3 \end{bmatrix},
  \quad
  Q = \sigma_{acc}\,I_3.
\]

\subsection{ARTVA Signal Model}

The buried victim's ARTVA transmitter is modelled as a single magnetic dipole
with moment $m$~\cite{azzollini2021es}, aligned with the world vertical axis.
The noiseless signal intensity at position $\mathbf{p}$ is:
\begin{equation}
  S(\mathbf{p}, \bm{\theta}) =
  m\,\frac{\sqrt{1 + 3\cos^2\!\vartheta}}{r^3},\quad
  r = \|\mathbf{p} - \bm{\theta}\|,
  \label{eq:artva}
\end{equation}
where $\bm{\theta}^*$ denotes the unknown source position,
$r$ is the Euclidean distance between the sensor and the source,
and $\cos\vartheta = (\mathbf{p} - \bm{\theta})_z / r$ is the cosine
of the angle between the relative position vector and the dipole axis.
Noisy measurements are modelled as
$\tilde{S}_i(t) = S(\mathbf{p}_i(t), \bm{\theta}^*) + \varepsilon_i(t)$,
with $\varepsilon_i \sim \mathcal{N}(0,\sigma^2)$.

In general, a buried victim may come to rest in any orientation, making
the dipole axis unknown. When the axis is not vertical, the signal
maximum shifts laterally from the victim's true position~\cite{azzollini2021es};
however, this is not a limitation specific to the proposed approach.
Any gradient-following strategy converges to the same signal maximum
$p^*$, which is precisely the point the drone marks for the rescuer to
begin excavation, i.e. the same reference a hand-held ARTVA search would
provide. The rescuer then refines the position empirically as digging
proceeds and the signal strengthens.


% ============================================================
\section{Proposed Solution}
\label{sec:solution}

\subsection{MPC Path-Following Controller}

Each drone tracks its mission waypoints via a receding-horizon MPC formulation.
At every step the IMDCL estimate $\hat{x}_i^+$ serves as initial condition and
the solver minimises:
\begin{equation}
\begin{split}
  \min_{u_0,\ldots,u_{N-1}} &\sum_{k=0}^{N-1}\!\bigl[
    w_p \|p_k - p_{\rm ref}\|^2 + w_a \|u_k\|^2\bigr]\\
  &+ w_{fp}\|p_N - p_{\rm ref}\|^2 + w_{fv}\|v_N\|^2
\end{split}
  \label{eq:mpc_cost}
\end{equation}
subject to \eqref{eq:dynamics} and box constraints
$\|v_j(k)\|_\infty \leq v_{\max} = 3\,$m/s,
$\|u_j(k)\|_\infty \leq a_{\max} = 6\,$m/s$^2$.
The problem is solved via the IPOPT interior-point solver through the CasADi
symbolic framework.

\paragraph{Warm-starting.}
At mission start a single initialisation solve runs up to $N_{\rm init} = 1000$
IPOPT iterations to obtain a high-quality primal--dual solution.
All subsequent MPC steps warm-start from the previous solution: the optimal
state trajectory is shifted by one step to initialise the primal variables.
IPOPT is then limited to $N_{\rm mpc} = 3$ iterations per step, keeping the
average solve time below $5\,$ms.

\paragraph{Waypoint sequencing.}
When the estimated distance to the active waypoint drops below
$\delta = 0.5\,$m, the mission planner advances to the next target and updates
$p_{\rm ref}$ without triggering a new initialisation solve.
The shifted warm-start from the current step carries over to the new reference,
ensuring smooth transitions between consecutive waypoints.

\subsection{ES: Extremum Seeking for Track Navigation}
\label{sec:es}

Extremum Seeking (ES) is a real-time, model-free optimisation technique that
finds the extremum of an unknown objective function by injecting a sinusoidal
perturbation into the search variable and demodulating the resulting output to
extract an approximate gradient, which is then used to drive a slow integrator
toward the optimum.  In this work the objective is the received ARTVA signal
strength and the search variable is the drone's horizontal position; the
perturbation produces a circular dither trajectory whose centre drifts
automatically toward the buried source without requiring any analytical model of
the magnetic field.

\begin{figure}[t]
  \centering
  \includegraphics[width=\linewidth]{ES_example.png}
  \caption{Illustration of the ES source-localization principle.
           The dither circle's centre drifts along the signal gradient,
           producing a spiralling trajectory that converges to the ARTVA
           source.  Figure adapted from~\cite{azzollini2021es}.}
  \label{fig:es_example}
\end{figure}

ES provides reactive gradient-following navigation that converges
asymptotically to the source, but yields a reliable position estimate only
after the dither amplitude has fully ramped up (${\approx}3\lambda \approx
45\,$s), so it used to guide the drone toward the source rather than estimate the source

On transition to \textsc{Track}, the drone switches to the gradient-free ES
controller.  The ES generates a continuously moving reference
$\mathbf{r}(t) = (x_{\rm ref}(t),\,y_{\rm ref}(t))$:
\begin{align}
  \dot{x}_{\rm ref} &= \sqrt{\alpha\,\omega}\cos(\omega t + \kappa\,y_t),
  \nonumber\\
  \dot{y}_{\rm ref} &= \sqrt{\alpha\,\omega}\sin(\omega t + \kappa\,y_t),
  \label{eq:es_ref}
\end{align}
where $y_t = S^{-1/3}$ is the conditioned ARTVA signal (convex, minimum at the
source), $\omega$ is the excitation frequency, $\kappa$ the feedback gain, and
$\alpha$ ramps from $0$ to $\alpha_{\max}$ via:
\begin{equation}
  \dot{\alpha} = \tfrac{1}{\lambda}(\alpha_{\max} - \alpha).
  \label{eq:es_ramp}
\end{equation}
The constraint $\sqrt{\alpha_{\max}\omega} = v_{\max} = 3\,$m/s ensures the
reference velocity never exceeds the actuator limit.  The trajectory is a circle
whose centre drifts along the signal gradient toward the source.  The MPC tracks
$\mathbf{r}(t)$ at each simulation step; DCGD runs concurrently on the
accumulated measurements.



\subsection{IMDCL: Decentralised Cooperative Localisation}

The Interim Master Decentralised Cooperative Localisation
(IMDCL)~\cite{kia2016cooperative} is a tightly coupled DCL algorithm that
replicates the result of a centralised EKF without requiring all-to-all
communication at every step.
The key insight is that cross-covariance terms between agents $i$ and $j$,
while never stored globally, can be reconstructed locally whenever a
relative measurement occurs.
To this end each agent $i$ maintains four quantities:
the state estimate $\hat{x}_i^+$, error covariance $P_i^+$, the cumulative
Jacobian product $\Phi_i$ (initialised to $I$ and multiplied by $F$ at every
propagation step), and local copies of the cross-covariance seed matrices
$\{\Pi^i_{jl}\}$ for all pairs $(j,l)$.
The cross-covariance between agents $i$ and $j$ at time $k$ can then be
recovered as $P_{ij}(k) = \Phi_i(k)\,\Pi^i_{ij}\,\Phi_j(k)^\top$ without
any additional communication.

\paragraph{Propagation (every step).}
Every agent independently propagates its local filter and advances $\Phi_i$:
\begin{align}
  \hat{x}_i^-(k{+}1) &= F\,\hat{x}_i^+(k) + G\,u_i(k), \nonumber\\
  P_i^-(k{+}1)       &= F\,P_i^+(k)\,F^\top + G\,Q\,G^\top, \nonumber\\
  \Phi_i(k{+}1)      &= F\,\Phi_i(k).  \label{eq:imdcl_prop}
\end{align}
No inter-agent communication is needed at this stage.

\paragraph{Update (when agent $a$ measures agent $b$).}
Agent $a$ is declared the \emph{interim master}.
It requests from $b$ the \emph{landmark-message}
$(\hat{x}_b^-,\,\Phi_b,\,P_b^-)$ and computes the whitened innovation
and innovation covariance:
\begin{align}
  r^a   &= z_{ab} - h_{ab}(\hat{x}_a^-,\hat{x}_b^-), \nonumber\\
  S_{ab} &= R^a
    + \tilde{H}_a P_a^- \tilde{H}_a^\top
    + \tilde{H}_b P_b^- \tilde{H}_b^\top \nonumber\\
    &\quad - \tilde{H}_a \Phi_a \Pi^a_{ab} \Phi_b^\top \tilde{H}_b^\top
    - (\cdot)^\top, \label{eq:imdcl_S}
\end{align}
where $\tilde{H}_a = [-I_3\;0_3]$, $\tilde{H}_b = [I_3\;0_3]$, and $z_{ab} =
p_b - p_a + \nu_{ab}$ is the UWB relative-position measurement.
The gain factors for the measuring agent $a$ and the landmark $b$ are
(Alg.~2, eq.~S9c in~\cite{kia2016cooperative}):
\begin{align}
  \Gamma_a &= \bigl(\Pi^a_{ab}\,\Phi_b^\top\tilde{H}_b^\top
              - \Phi_a^{-1}P_a^-\tilde{H}_a^\top\bigr)S_{ab}^{-1/2},
              \label{eq:gamma_a}\\
  \Gamma_b &= \bigl(\Phi_b^{-1}P_b^-\tilde{H}_b^\top
              - \Pi^b_{ba}\,\Phi_a^\top\tilde{H}_a^\top\bigr)S_{ab}^{-1/2}.
              \label{eq:gamma_b}
\end{align}
Agent $a$ packs
$\bigl(\Gamma_a,\,\Gamma_b,\;\Phi_b^\top\tilde{H}_b^\top S_{ab}^{-1/2},\;
\Phi_a^\top\tilde{H}_a^\top S_{ab}^{-1/2}\bigr)$
into the \emph{update-message} and broadcasts it to the whole team.
Every agent $i\notin\{a,b\}$ derives its own gain locally from its $\Pi$
copies (eq.~S11 in~\cite{kia2016cooperative}):
\begin{equation}
  \Gamma_i = \Pi^i_{ib}\,\Phi_b^\top\tilde{H}_b^\top S_{ab}^{-1/2}
           - \Pi^i_{ia}\,\Phi_a^\top\tilde{H}_a^\top S_{ab}^{-1/2},
  \label{eq:gamma_i}
\end{equation}
and every agent $i \in \mathcal{V}$ updates:
\begin{align}
  \hat{x}_i^+ &= \hat{x}_i^- + \Phi_i\,\Gamma_i\,\tilde{r}^a, \label{eq:imdcl_xup}\\
  P_i^+        &= P_i^- - \Phi_i\,\Gamma_i\,\Gamma_i^\top\,\Phi_i^\top, \label{eq:imdcl_Pup}\\
  \Pi^i_{jl}(k{+}1) &= \Pi^i_{jl}(k) - \Gamma_j\,\Gamma_l^\top, \label{eq:imdcl_Piup}
\end{align}
where $\tilde{r}^a = S_{ab}^{-1/2}r^a$.
The seed update \eqref{eq:imdcl_Piup} propagates the measurement benefit
to all cross-covariance pairs, and $\Phi_i$ is reset to $I$ after the update.

\paragraph{Measurement sources.}
Two exteroceptive sensors drive the filter updates:
\begin{itemize}
  \item \textbf{LiDAR altimeter}: absolute altitude above terrain,
        processed as an absolute update with $\sigma_{\rm LiDAR} = 0.05\,$m.
  \item \textbf{UWB inter-drone ranging}: 3-D relative position between
        all pairs within $R_c$, processed as a relative update via
        \eqref{eq:imdcl_S}--\eqref{eq:imdcl_Piup} with $\sigma_{\rm rel} =
        0.3\,$m.
\end{itemize}

\subsection{DCGD: Distributed Source Estimation}

While the extremum-seeking controller drives the Track drones towards
the source, in this work it is not used to produce an explicit position estimate.
For that purpose, the same drones cooperatively minimise:
\begin{equation}
  J(\bm{\theta}) = \sum_{i \in \mathcal{T}}
    \sum_{(\mathbf{p},\tilde{s}) \in \mathcal{B}_i}
    \bigl(\tilde{s} - S(\mathbf{p}, \bm{\theta})\bigr)^2,
\end{equation}
where $\mathcal{B}_i$ is a sliding batch of $|\mathcal{B}_i| = 5$ measurements
and $\mathcal{T}$ is the set of active Track drones.  Each DCGD iteration
consists of two steps.

\paragraph{Adapt.}
Each agent computes a normalised gradient descent step via a centred
finite-difference gradient ($h = 0.1\,$m) with step size $\alpha = 0.3\,$m:
\begin{equation}
  \bm{\theta}_i^{(k+\frac{1}{2})} = \bm{\theta}_i^{(k)}
  - \alpha\,\frac{\bar{g}_i(\bm{\theta}_i^{(k)})}
                 {\|\bar{g}_i(\bm{\theta}_i^{(k)})\|}.
  \label{eq:adapt}
\end{equation}
Normalisation enforces a fixed step size $\alpha$ in parameter space
regardless of signal magnitude, which is critical given the $1/r^3$
variation of the ARTVA field as drones approach the source.
It also avoids the use of adaptive optimisers (e.g.\ Adam, RMSprop), whose
per-coordinate gradient statistics would require additional
inter-drone communication beyond the single Combine step.

\paragraph{Combine.}
A single weighted average consensus step fuses the local estimates:
\begin{equation}
  \bm{\theta}_i^{(k+1)} = (1-\beta)\,\bm{\theta}_i^{(k+\frac{1}{2})}
  + \frac{\beta}{|\mathcal{N}_i \cap \mathcal{T}|}
    \sum_{j \in \mathcal{N}_i \cap \mathcal{T}}
    \bm{\theta}_j^{(k+\frac{1}{2})},
  \label{eq:combine}
\end{equation}
with $\beta = 0.4$.  The self-weight $(1-\beta) = 0.6$ keeps the local
gradient information from the Adapt step dominant, while $\beta$ is large
enough to propagate consensus across the network within the communication
radius $R_c$.  This implements the ATC diffusion
strategy~\cite{chen2012diffusion}.

\paragraph{Spatial diversity and support orbits.}
The accuracy of the cooperative source estimate depends critically on the
\emph{spatial diversity} of the measurement set: if all Track drones
converge to nearly the same position, their observations become highly
correlated and DCGD effectively reduces to a single-agent estimator,
leaving the nonlinear ARTVA field poorly identified.
To counteract this, the two Support drones recruited via min-consensus are
assigned circular orbits of opposite handedness around the Stop drone, one
clockwise, one counter-clockwise, at the estimated source distance.
This anti-phase orbit assignment guarantees that, at every instant, the three
active drones span a wide angular sector around the source, maximising the
geometric diversity of the measurement footprint and thereby accelerating
DCGD convergence toward an accurate triangulation.



\subsection{Min-Consensus Partner Selection}

When the first Track drone $h$ reaches \textsc{Stop} (strong-signal threshold),
the system selects the two nearest reachable drones as support partners fully
distributedly.

\paragraph{Step 0 — Connected component.}
A BFS (\textit{Breadth-First Search}) from $h$ identifies $\mathcal{C}(h)$, the set of drones reachable through
$\mathcal{G}$.  Drones in $\mathcal{V} \setminus \mathcal{C}(h)$ are excluded.

\paragraph{Step 1 — Initialisation.}
Each $i \in \mathcal{C}(h)$ maintains a distance row
$\mathbf{D}_i \in (\mathbb{R} \cup \{+\infty\})^{|\mathcal{C}(h)|}$,
initialised to $+\infty$ everywhere except $\mathbf{D}_h[h] = 0$.

\paragraph{Step 2 — Iterations.}
For up to $K_{\max} = 10$ rounds two updates propagate synchronously:
(i)~\emph{awareness}: newly-informed drones set $\mathbf{D}_i[i] =
\|\mathbf{p}_i - \mathbf{p}_h\|$;
(ii)~\emph{min-update}: $\mathbf{D}_i \leftarrow \min(\mathbf{D}_i,
\mathbf{D}_j)$ for each neighbour $j$.  Convergence is guaranteed in at most
$\mathrm{diam}(\mathcal{G}[\mathcal{C}(h)])$ rounds.

\paragraph{Step 3 — Selection.}
The agreed distance for candidate $c_k$ is
$d^*_k = \min_{i \in \mathcal{C}(h)} \mathbf{D}_i[k]$.
The $P$ partners with smallest $d^*_k$ are selected.

% ============================================================
\section{Implementation}
\label{sec:impl}

\subsection{Noise Calibration}

Before flight each drone collects $M = 20$ ARTVA background measurements and
estimates its local noise standard deviation $\hat{\sigma}_i$.  Ten iterations of
average consensus~\cite{olfati2004consensus} drive all agents to a common estimate
$\hat{\sigma}$, which sets:
\begin{itemize}
  \item \textbf{Detection threshold}:
        $\tau_{\rm det} = c_{\rm det}\,\hat{\sigma}$ (SEARCH $\to$ TRACK),
        with $c_{\rm det} \in \{50, 100\}$, floored at the signal strength
        corresponding to $R = 50\,$m from the source (ES convergence range);
  \item \textbf{Stop threshold}:
        $\tau_{\rm stop} = 1000\,\hat{\sigma}$ (TRACK/SUPPORT $\to$ STOP).
\end{itemize}
For the nominal $\sigma = 10^{-6}$ and $c_{\rm det} = 100$, these correspond to
detection at $\sim\!20\,$m and stopping at $\sim\!5\,$m from the source.  Both
thresholds scale automatically with the measured noise.

\subsection{Mission State Machine}

Each drone executes a four-state finite automaton (Fig.~\ref{fig:fsm}):
\begin{itemize}
  \item \textsc{Search}: boustrophedon lawnmower coverage of an assigned
        workspace lane; the area is partitioned equally along the $x$-axis.
  \item \textsc{Track}: the ES controller generates a continuously drifting
        circular reference converging toward the ARTVA source; DCGD runs
        concurrently.
  \item \textsc{Stop}: strong-signal hovering; triggers partner selection and
        orbit assignment.
  \item \textsc{Support}: circular orbit (CW or CCW) around the Stop drone at
        the estimated source distance, providing additional measurement diversity
        for DCGD.
\end{itemize}
\textsc{Search} and \textsc{Support} transitions are arrival-gated: waypoint
updates are issued only when the current target is within $\delta = 0.3\,$m of
the estimated drone position.  During \textsc{Track} the ES reference is updated
every simulation step without an arrival gate.

\begin{figure}[t]
  \centering
  \framebox{\parbox{0.9\linewidth}{\centering\vspace{1.2cm}
    \textit{[Mission state machine diagram]}
  \vspace{1.2cm}}}
  \caption{Mission state machine for each drone agent.  Arrows indicate
           transitions with their respective trigger conditions.}
  \label{fig:fsm}
\end{figure}

\subsection{Refinement Phase}

When $\geq 3$ drones reach \textsc{Stop}, DCGD executes an additional
$K_{\rm refine} = 50$ iterations without drone movement.  This stationary
refinement exploits the measurement batch accumulated near the source to
sharpen the estimate before reporting.

\subsection{Parameter Summary}

Table~\ref{tab:params} collects all simulation and algorithm parameters.

\begin{table}[t]
\centering
\caption{Simulation and algorithm parameters.}
\label{tab:params}
\begin{tabular}{lll}
\toprule
Parameter & Symbol / key & Value \\
\midrule
\multicolumn{3}{l}{\textit{MPC controller}} \\
Horizon              & $N$         & 20 steps \\
Sampling time        & $dt$        & 0.1\,s   \\
Max.\ acceleration   & $a_{\max}$  & 6\,m/s$^2$ \\
Max.\ velocity       & $v_{\max}$  & 3\,m/s   \\
Pos.\ weight         & $w_p$       & 100      \\
Input weight         & $w_a$       & 0.1      \\
Terminal pos.\ wt.   & $w_{fp}$    & 1\,000   \\
Terminal vel.\ wt.   & $w_{fv}$    & 10       \\
\midrule
\multicolumn{3}{l}{\textit{IMDCL filter}} \\
Initial pos.\ std    & $\sqrt{P_0^{pp}}$ & 0.5\,m   \\
Relative meas.\ std  & $\sigma_{\rm rel}$ & 0.3\,m  \\
LiDAR meas.\ std     & $\sigma_{\rm LiDAR}$ & 0.05\,m \\
Process noise        & $\sigma_{acc}$ & 0.05\,m/s$^2$ \\
Comm.\ radius        & $R_c$          & 80\,m     \\
\midrule
\multicolumn{3}{l}{\textit{Extremum Seeking (Track)}} \\
Max.\ amplitude      & $\alpha_{\max}$ & 20      \\
Frequency            & $\omega$        & 0.45\,rad/s \\
Feedback gain        & $\kappa$        & 0.05    \\
Ramp time const.     & $\lambda$       & 15\,s   \\
\midrule
\multicolumn{3}{l}{\textit{DCGD source estimator}} \\
Step size            & $\alpha$        & 0.3\,m  \\
Consensus weight     & $\beta$         & 0.4     \\
Finite-diff.\ step   & $h$             & 0.1\,m  \\
Measurement batch    & $|\mathcal{B}_i|$ & 5     \\
Refinement iters     & $K_{\rm refine}$ & 300    \\
\midrule
\multicolumn{3}{l}{\textit{Partner selection}} \\
Partners             & $P$            & 2       \\
Max.\ consensus rounds & $K_{\max}$   & 10      \\
\midrule
\multicolumn{3}{l}{\textit{Noise calibration}} \\
Calibration samples  & $M$            & 20      \\
Detection factor     & $c_{\rm det}$  & 50 or 100 \\
Stop factor          & $c_{\rm stop}$ & 1\,000  \\
\midrule
\multicolumn{3}{l}{\textit{ARTVA model}} \\
Dipole moment        & $m$            & 1.0     \\
Nominal noise std    & $\sigma$       & $10^{-6}$ \\
\bottomrule
\end{tabular}
\end{table}

% ============================================================
\section{Experimental Results}
\label{sec:results}

\subsection{Individual Algorithm Validation}

\subsubsection{MPC Path-Following Performance}

The warm-start strategy keeps the average IPOPT solve time below $5\,$ms per step,
using only three interior-point iterations after the initial full solve.
Steady-state waypoint tracking error remains below $0.15\,$m across all runs,
well within the arrival threshold $\delta = 0.3\,$m.
Fig.~\ref{fig:mpc_traj} shows representative 3-D trajectories and control-input
norms for a three-drone team covering a $200\times200\,$m lawnmower pattern.

\begin{figure}[t]
  \centering
  \framebox{\parbox{0.9\linewidth}{\centering\vspace{1.5cm}
    \textit{[MPC trajectories and acceleration norm vs.\ time]}
  \vspace{1.5cm}}}
  \caption{MPC trajectories (top) and control input magnitude (bottom) for three
           drones in the \textsc{Search} phase.  Dashed lines indicate waypoint
           targets; grey line marks $a_{\max} = 6\,$m/s$^2$.}
  \label{fig:mpc_traj}
\end{figure}

\subsubsection{Cooperative Localisation}

Fig.~\ref{fig:imdcl_error} shows position estimation error for a representative
drone over a 500-step GPS-denied simulation with $N = 5$.
The IMDCL filter combines UWB relative measurements ($\sigma_{\rm rel} = 0.3\,$m)
and LiDAR altitude corrections ($\sigma_{\rm LiDAR} = 0.05\,$m) to maintain
bounded RMS error below $0.5\,$m.
Cross-covariance correction prevents the information double-counting that would
otherwise inflate filter confidence in a naïve distributed EKF.

\begin{figure}[t]
  \centering
  \framebox{\parbox{0.9\linewidth}{\centering\vspace{1.5cm}
    \textit{[IMDCL position error vs.\ time]}
  \vspace{1.5cm}}}
  \caption{Position estimation error (IMDCL) for one drone over 500 steps.
           Shaded band: $\pm 1\sigma$ across 20 Monte Carlo runs.}
  \label{fig:imdcl_error}
\end{figure}

\subsubsection{DCGD Source Estimation}

Fig.~\ref{fig:dcgd_conv} reports DCGD convergence during the Track phase.
The source estimate error decreases monotonically, reaching sub-metre accuracy
within 150 iterations for $N_{\mathcal{T}} = 3$ Track drones.  The ATC Combine
step consistently accelerates convergence relative to isolated local gradient
descent, confirming the benefit of the diffusion architecture.

\begin{figure}[t]
  \centering
  \framebox{\parbox{0.9\linewidth}{\centering\vspace{1.5cm}
    \textit{[DCGD source estimation error vs.\ iteration]}
  \vspace{1.5cm}}}
  \caption{Source estimation error $\|\bm{\theta}_i - \bm{\theta}^*\|$ for three
           Track drones.  Dashed: local gradient descent only; solid: DCGD with
           ATC Combine.}
  \label{fig:dcgd_conv}
\end{figure}

\subsubsection{Partner Selection}

Table~\ref{tab:consensus} summarises min-consensus partner selection results for
three representative graph topologies.  In all cases the algorithm converges in at
most 4 rounds (well within $K_{\max} = 10$) and selects the geometrically correct
partners.

\begin{table}[t]
\centering
\caption{Min-consensus partner selection: rounds to convergence and selected
         partners for $N = 6$, $P = 2$.}
\label{tab:consensus}
\begin{tabular}{cccc}
\toprule
Scenario & $|\mathcal{C}(h)|$ & Rounds & Partners \\
\midrule
Dense graph   & 6 & 2 & $\{2, 3\}$ \\
Sparse graph  & 5 & 4 & $\{3, 5\}$ \\
Split network & 3 & 2 & $\{2, 3\}$ \\
\bottomrule
\end{tabular}
\end{table}

\subsection{Parametric Sweep}
\label{sec:sweep}

\subsubsection{Experimental Setup}

Five workspace patches of side $A \in \{100, 200\}\,$m are extracted from the
real TINItaly 10\,m-resolution DEM (tile \texttt{w51065\_s10.tif}, Trentino Alps)
at centres distributed across the tile, yielding genuinely distinct alpine terrain
profiles: ridgelines, bowls, and gentle slopes (Fig.~\ref{fig:dem_workspaces}).
For each parameter combination, victim positions are sampled uniformly inside
$[0.10,\,0.90]^2$ of the workspace (relative coordinates) and burial depth
uniformly in $[1.0,\,5.0]\,$m.  The full sweep grid is shown in
Table~\ref{tab:sweep}; results are recorded in a CSV file with automatic resume
after interruption.

\begin{figure}[t]
  \centering
  \framebox{\parbox{0.9\linewidth}{\centering\vspace{1.5cm}
    \textit{[TINItaly DEM with five selected workspace patches]}
  \vspace{1.5cm}}}
  \caption{TINItaly DEM with the five workspace patches used in the parametric
           sweep.  Each coloured rectangle corresponds to one terrain configuration.}
  \label{fig:dem_workspaces}
\end{figure}

\begin{table}[t]
\centering
\caption{Parametric sweep axes and ranges.}
\label{tab:sweep}
\begin{tabular}{ll}
\toprule
Parameter & Values \\
\midrule
Workspace side $A$ [m] & 100, 200 \\
Fleet size $N$         & 3, 4, 5 \\
Victim positions       & [XX] per combination (random) \\
Burial depth [m]       & $\mathcal{U}(1.0,\,5.0)$ \\
ARTVA noise $\sigma$   & $10^{-7}$, $10^{-6}$, $10^{-5}$ \\
Detection factor $c_{\rm det}$ & 50, 100 \\
Comm.\ radius $R_c$ [m] & 25, 50, 80, 120 \\
Terrain patches         & 5 \\
\midrule
\textbf{Total runs}    & \textbf{[XXXX]} \\
\bottomrule
\end{tabular}
\end{table}

Each run records: (i)~\emph{found} (Boolean: $\geq 3$ drones in \textsc{Stop}
within 15 min and 2-D estimate error $<10\,$m); (ii)~\emph{time\_found} [s];
(iii)~\emph{pos\_std} [m], the square root of the summed per-axis variance of
individual drone estimates (inter-drone consensus quality);
(iv)~\emph{est\_error\_2d} [m], the Euclidean error between the mean estimate and
the true victim XY position.

\subsubsection{Overall Performance}

Across all [XXXX] runs the system achieves an overall mission success rate of
\textbf{[XX.X]\,\%} ([XXXX] out of [XXXX] runs).
Table~\ref{tab:results} summarises key metrics broken down by the main sweep axes.
Fig.~\ref{fig:heatmap} shows the success-rate heatmap as a function of noise and
communication radius.

\begin{table}[t]
\centering
\caption{Sweep results broken down by parameter.}
\label{tab:results}
\begin{tabular}{llcc}
\toprule
Axis & Value & Success [\%] & Median time [s] \\
\midrule
\multirow{3}{*}{Fleet $N$}
  & 3 & [XX.X] & [XX] \\
  & 4 & [XX.X] & [XX] \\
  & 5 & [XX.X] & [XX] \\
\midrule
\multirow{2}{*}{Area $A$}
  & 100\,m & [XX.X] & --- \\
  & 200\,m & [XX.X] & --- \\
\midrule
\multirow{3}{*}{Noise $\sigma$}
  & $10^{-7}$ & [XX.X] & --- \\
  & $10^{-6}$ & [XX.X] & --- \\
  & $10^{-5}$ & [XX.X] & --- \\
\midrule
\multirow{4}{*}{Radius $R_c$}
  & 25\,m  & [XX.X] & --- \\
  & 50\,m  & [XX.X] & --- \\
  & 80\,m  & [XX.X] & --- \\
  & 120\,m & [XX.X] & --- \\
\bottomrule
\end{tabular}
\end{table}

\begin{figure}[t]
  \centering
  \framebox{\parbox{0.9\linewidth}{\centering\vspace{1.5cm}
    \textit{[Success-rate heatmap: noise $\times$ comm.\ radius, by fleet size]}
  \vspace{1.5cm}}}
  \caption{Success rate [\%] as a function of ARTVA noise $\sigma$ and
           communication radius $R_c$ for each fleet size and workspace area.}
  \label{fig:heatmap}
\end{figure}

\subsubsection{Effect of Fleet Size and Workspace Area}

Larger fleet sizes consistently improve both success rate and search time.
Moving from $N = 3$ to $N = 5$ raises success by $\sim[X.X]$\,pp and reduces
median search time by $\sim[X]$\,s, owing to faster lawnmower coverage and more
robust DCGD consensus.  The 200\,m workspace is harder: larger area increases the
distance to the victim and raises the risk of the fleet never reaching the Track
phase within the 15-min timeout.  Fig.~\ref{fig:time_boxplot} shows the
distribution of search times for successful runs.

\begin{figure}[t]
  \centering
  \framebox{\parbox{0.9\linewidth}{\centering\vspace{1.5cm}
    \textit{[Boxplot of search time by fleet size]}
  \vspace{1.5cm}}}
  \caption{Distribution of search time (successful runs only) for each fleet size
           and workspace area.  Annotated medians in seconds.}
  \label{fig:time_boxplot}
\end{figure}

\subsubsection{Effect of Communication Radius}

Communication radius critically governs graph connectivity and partner selection
reliability.  At $R_c = 25\,$m the network is frequently fragmented during the
lawnmower phase, preventing the min-consensus algorithm from reaching the $P = 2$
required partners and leaving Track drones without support orbits.  At
$R_c = 120\,$m success reaches [XX.X]\,\%.  Fig.~\ref{fig:comm_radius} shows the
monotonic improvement with $R_c$ for all fleet sizes.

\begin{figure}[t]
  \centering
  \framebox{\parbox{0.9\linewidth}{\centering\vspace{1.5cm}
    \textit{[Success rate vs.\ comm.\ radius for each fleet size]}
  \vspace{1.5cm}}}
  \caption{Success rate as a function of UWB communication radius $R_c$ for each
           fleet size (left: 100\,m workspace; right: 200\,m workspace).}
  \label{fig:comm_radius}
\end{figure}

\subsubsection{Effect of Noise Level}

ARTVA noise has the strongest single-factor impact on mission success.
At $\sigma = 10^{-5}$ (highest noise tested) the detection threshold is crossed
only at very short range, leaving insufficient time for ES to converge before the
timeout.  At $\sigma = 10^{-7}$ the system operates in near-ideal conditions and
consistently localises the victim within [XX]\,s.  The automatic noise-calibration
procedure scales both thresholds proportionally, so the [XX.X]\,pp gap across
noise levels primarily reflects the reduced ES convergence range rather than
threshold miscalibration.

\subsubsection{Source Localisation Accuracy}

The 2-D source localisation error (mean [XX.X]\,m, median [XX.X]\,m over all
successful runs) reflects the ES convergence characteristic: the \textsc{Stop}
transition fires when the signal threshold is crossed, which corresponds to being
within a few metres of the source in 3-D but can project to a larger 2-D footprint
depending on burial depth and terrain slope.  The subsequent 300-iteration
stationary DCGD refinement sharpens the estimate further, but residual uncertainty
is dominated by the burial-depth ambiguity inherent in the planar 2-D metric.
Fig.~\ref{fig:cdf} shows the empirical CDF of search time for each fleet size.

\begin{figure}[t]
  \centering
  \framebox{\parbox{0.9\linewidth}{\centering\vspace{1.5cm}
    \textit{[CDF of search time by fleet size and area]}
  \vspace{1.5cm}}}
  \caption{Empirical CDF of search time for successful runs.  Dashed horizontal
           line at 90\,\%.}
  \label{fig:cdf}
\end{figure}

% ============================================================
\section{Conclusions}
\label{sec:conclusions}

This paper presented a complete distributed multi-drone system for ARTVA victim
localisation, integrating MPC path-following, IMDCL cooperative localisation, ES
source tracking, DCGD source estimation, and min-consensus partner selection.
All components operate with only local state and within a limited communication
radius, requiring no central coordinator.  Raw measurements never leave the drone:
only compact event messages are exchanged.

A parametric sweep of [XXXX] experiments spanning five real TINItaly alpine
terrain patches, three fleet sizes, three noise levels, four communication radii,
and randomised burial depths achieves an overall mission success rate of
[XX.X]\,\% with a median search time of [XX]\,s.  Fleet size has the most
consistent positive effect on both success and speed, while communication radius
critically governs connectivity and partner selection reliability.  ARTVA noise
has the strongest negative impact, as it shrinks the ES convergence window.

Key limitations are: (i)~the 2-D localisation error ($\sim[XX]\,$m median) driven
by burial-depth ambiguity and the ES spiral stop condition; (ii)~sensitivity of
ES convergence to SNR at the detection boundary; and (iii)~the single buried
source assumption.

Future directions include multi-source extension via spatial clustering of
measurement batches, tighter 3-D localisation by fusing DCGD with bearing
information from multiple orbit drones, replacement of finite-difference gradients
with analytic Jacobians of \eqref{eq:artva}, and experimental validation on
physical drone platforms over real snowpack.

% ============================================================
\begin{thebibliography}{9}

\bibitem{azzollini2021es}
  I.~A.~Azzollini, L.~Gentilini, A.~Bassi, E.~Bernuau, P.~Colaneri,
  E.~D'Amato, and S.~Fabbri,
  ``Extremum seeking based on the gradient control: Design and analysis,''
  \textit{arXiv preprint arXiv:2106.14514}, 2021.

\bibitem{kia2016cooperative}
  S.~S.~Kia, S.~Rounds, and S.~Mart\'inez,
  ``Cooperative localization for mobile agents: A recursive decentralized
    algorithm based on Kalman-filter decentralization,''
  \textit{IEEE Control Syst.\ Mag.}, vol.~36, no.~2, pp.~86--101, 2016.

\bibitem{olfati2004consensus}
  R.~Olfati-Saber and R.~M.~Murray,
  ``Consensus problems in networks of agents with switching topology and
    time-delays,''
  \textit{IEEE Trans.\ Autom.\ Control}, vol.~49, no.~9,
  pp.~1520--1533, 2004.

\bibitem{chen2012diffusion}
  J.~Chen and A.~H.~Sayed,
  ``Diffusion adaptation strategies for distributed optimization and learning
    over networks,''
  \textit{IEEE Trans.\ Signal Process.}, vol.~60, no.~8,
  pp.~4289--4305, 2012.

\bibitem{artva2020}
  M.~Romanoni, F.~Furlan, and M.~Matteucci,
  ``Autonomous ARTVA-based search with a UAV,''
  in \textit{Proc.\ IEEE/RSJ IROS}, Las Vegas, NV, USA, 2020.

\bibitem{rawlings2017mpc}
  J.~B.~Rawlings, D.~Q.~Mayne, and M.~Diehl,
  \textit{Model Predictive Control: Theory, Computation, and Design},
  2nd~ed.\ Nob Hill Publishing, Santa Barbara, CA, USA, 2017.

\end{thebibliography}

\end{document}
